%% Fasit til kortsett fra mattelab.no – laget automatisk av bygg.py.
%% Kompiler med lualatex.
\documentclass[11pt,a4paper]{article}
\usepackage[norsk]{babel}
\usepackage[a4paper,margin=1.2cm,bottom=1.7cm]{geometry}
\usepackage{fontspec}
\usepackage{unicode-math}
\setmainfont{TeX Gyre Pagella}
\setsansfont{TeX Gyre Heros}[Scale=MatchLowercase]
\setmathfont{TeX Gyre Pagella Math}
\usepackage{xcolor}
\definecolor{blekk}{HTML}{16213A}
\definecolor{dempet}{HTML}{566077}
\definecolor{m1}{HTML}{2F4AB0}
\usepackage{tikz}
\usepackage{enumitem}
\usepackage{qrcode}
\usepackage{fancyhdr}
\usepackage[hidelinks]{hyperref}
\pdfvariable suppressoptionalinfo \numexpr 1+2+4+8+16+32+64+128+256+512\relax
\setlength{\parindent}{0pt}
\setlength{\parskip}{0pt}
\pagestyle{fancy}
\fancyhf{}
\renewcommand{\headrulewidth}{0pt}
\fancyfoot[L]{\sffamily\footnotesize\color{dempet}mattelab.no · Definisjonen av den deriverte · kortsett K12 · fasit}
\fancyfoot[R]{\sffamily\footnotesize\color{dempet}Side \thepage}
\begin{document}
\begin{minipage}[t]{0.8\linewidth}\vspace{0pt}
{\sffamily\small\color{dempet}Fasit for læreren · kortsett K12 · R1}\par\vspace{0.25em}
{\fontsize{20}{24}\selectfont\bfseries Definisjonen av den deriverte}\par\vspace{0.5em}
Finn fire og fire kort som hører sammen: en funksjon f, differansen f(x + h) − f(x), differansekvotienten forenklet, og den deriverte. Regn selv, og kontroller at kortene stemmer.\par\vspace{0.4em}{\small\color{dempet}Kortene er nummerert i den rekkefølgen de står i utskriften.}
\end{minipage}\hfill
\begin{minipage}[t]{0.15\linewidth}\vspace{0pt}\raggedleft
\qrcode[height=1.9cm]{https://mattelab.no/varierende/definisjonen-av-den-deriverte.html}
\end{minipage}\par\vspace{0.8em}

\begin{enumerate}[leftmargin=1.6em, itemsep=0.7em]
\item \begin{minipage}[t]{\linewidth}\textsf{\small K12-15}\quad $f'(x) = 2x$ \\ \textsf{\small K12-21}\quad $f(x) = x^2$ \\ \textsf{\small K12-22}\quad $f(x + h) - f(x) = 2xh + h^2$ \\ \textsf{\small K12-24}\quad $\dfrac{f(x + h) - f(x)}{h} = 2x + h$\end{minipage}
\item \begin{minipage}[t]{\linewidth}\textsf{\small K12-11}\quad $\dfrac{f(x + h) - f(x)}{h} = 4x + 2h$ \\ \textsf{\small K12-14}\quad $f'(x) = 4x$ \\ \textsf{\small K12-16}\quad $f(x) = 2x^2$ \\ \textsf{\small K12-23}\quad $f(x + h) - f(x) = 4xh + 2h^2$\end{minipage}
\item \begin{minipage}[t]{\linewidth}\textsf{\small K12-09}\quad $f(x + h) - f(x) = 2xh + h^2 + 3h$ \\ \textsf{\small K12-10}\quad $f'(x) = 2x + 3$ \\ \textsf{\small K12-12}\quad $f(x) = x^2 + 3x$ \\ \textsf{\small K12-17}\quad $\dfrac{f(x + h) - f(x)}{h} = 2x + h + 3$\end{minipage}
\item \begin{minipage}[t]{\linewidth}\textsf{\small K12-01}\quad $f(x) = 3x - 1$ \\ \textsf{\small K12-07}\quad $f(x + h) - f(x) = 3h$ \\ \textsf{\small K12-08}\quad $\dfrac{f(x + h) - f(x)}{h} = 3$ \\ \textsf{\small K12-18}\quad $f'(x) = 3$\end{minipage}
\item \begin{minipage}[t]{\linewidth}\textsf{\small K12-03}\quad $f(x) = x^3$ \\ \textsf{\small K12-04}\quad $\dfrac{f(x + h) - f(x)}{h} = 3x^2 + 3xh + h^2$ \\ \textsf{\small K12-06}\quad $f'(x) = 3x^2$ \\ \textsf{\small K12-19}\quad $f(x + h) - f(x) = 3x^2h + 3xh^2 + h^3$\end{minipage}
\item \begin{minipage}[t]{\linewidth}\textsf{\small K12-02}\quad $f'(x) = -\dfrac{1}{x^2}$ \\ \textsf{\small K12-05}\quad $f(x + h) - f(x) = -\dfrac{h}{x(x + h)}$ \\ \textsf{\small K12-13}\quad $\dfrac{f(x + h) - f(x)}{h} = -\dfrac{1}{x(x + h)}$ \\ \textsf{\small K12-20}\quad $f(x) = \dfrac{1}{x}$\end{minipage}
\end{enumerate}
\subsection*{Til læreren}
Fellene: $x^2$ og $2x^2$ gir uttrykk som ligner, og det samme gjør $x^2$ og $x^2 + 3x$. For $3x - 1$ forsvinner h helt, så den deriverte er konstant. Spør: Hvorfor kan vi ikke sette h = 0 i differansekvotienten før vi har forkortet? Hvilket ledd i differansen forsvinner alltid når vi trekker fra f(x)?
\end{document}
